
\section*{Supplementary Information}

**A benchmark cannot order more methods than its labels can distinguish:
capacity limits in CAID, and the protocol that escapes them**

Tommaso R. Marena


\vspace{0.5em}\hrule\vspace{0.8em}



\section*{Contents}

\begin{center}\footnotesize
\begin{tabular}{ll}
\toprule
\textbf{} & \textbf{} \\
\midrule
\textbf{Note S1} & The capacity theorem: statement, proof outline, and scope \\
\textbf{Note S2} & The second-order identity, its hypotheses, and two corrections \\
\textbf{Note S3} & The protocol, worked, including what it refuses to order \\
\textbf{Note S4} & Pre-registered variants and their outcomes --- five negatives \\
\textbf{Note S5} & Errors found and corrected, with the measurement that caught each \\
\textbf{Note S6} & Inventory of the formal development \\
\textbf{Note S7} & The screen: arbitrary dependence, and the unit of testing \\
\textbf{Table S1} & AUC decomposition, every method $\times$ every reference (CSV, 430 rows) \\
\textbf{Table S2} & Inversion certificates against the pooled winner (CSV) \\
\textbf{Table S3} & Pairwise leaderboard, complete, five references (CSV, 357 rows) \\
\textbf{Table S4} & Operating cost, whole field, three tolerated miss rates (CSV, 1,290 rows) \\
\textbf{Table S5} & All thirty per-target paired tests with Holm adjustment (CSV) \\
\textbf{Table S6} & Capacity by benchmark, three independent routes (CSV) \\
\textbf{Table S7} & Per-protein label and pairwise noise from MobiDB (CSV, 159 rows) \\
\textbf{Table S8} & Cross-round reproducibility, all protocol pairings (CSV) \\
\textbf{Table S9} & Coverage as a confound --- what declining targets is worth \\
\textbf{Table S10} & Thresholded error counts at each method's own operating point \\
\textbf{Table S11} & Certification frontier against assumed annotation error \\
\textbf{Table S12} & Data provenance: md5s and composition \\
\textbf{Table S13} & Job registry \\
\textbf{Table S14} & Registered endpoints closed for this paper (CSV) \\
\textbf{Table S15} & The region-screen instantiation, both rounds (CSV) \\
\textbf{Table S16} & The screen as run: discoveries, FDP, power, calibration budget (CSV) \\
\bottomrule
\end{tabular}
\end{center}

CSV tables are in \code{paper/supplementary/} and are regenerated by
\code{paper/make\_supplementary.py} directly from the analysis outputs.


\vspace{0.5em}\hrule\vspace{0.8em}



\section*{Note S1 --- The capacity theorem}


\subsection*{Statement}

Let a benchmark score \code{n} items, each method receiving a score that is an
average over those items and therefore lies on the grid of denominator \code{n}. Let
the annotation be wrong on at most a fraction \code{ε} of items (a worst-case budget,
not a probability model), and let \code{δ ≥ 0} be the smallest score difference worth
calling a difference. Write \code{c = max(δ, 2ε)}. Then any family of methods that
the benchmark places in a \emph{certified} total order has cardinality at most

\begin{quote}\ttfamily\footnotesize\begin{flushleft}
k(n,~ε,~δ)~=~n~/~(⌊c·n⌋~+~1)~+~1.\\
\end{flushleft}\end{quote}

\code{BenchmarkCapacity.card\_le\_benchCapacity}, Lean 4, sorry-free, axioms
\code{propext}, \code{Classical.choice}, \code{Quot.sound}.

\textbf{Sharpness.} \code{benchCapacity\_attained} constructs a family of exactly \code{k}
methods that the benchmark does resolve, so \code{k} is the capacity and not merely
an upper bound. The theorem and its attainment instance are given in three score
models --- values on the grid \code{\{0, 1/n, …, 1\}}, integer error counts in \code{\{0,…,N\}},
and real scores in \code{[0, R]} --- because the first is what CAID reports, the second
is what a thresholded metric reports, and the third is what a referee will ask
about.

\textbf{The n-free corollary.} \code{benchCapacity\_noise\_only}: with \code{δ = 0},

\begin{quote}\ttfamily\footnotesize\begin{flushleft}
k~≤~⌈1/(2ε)⌉~~~~for~every~n.\\
\end{flushleft}\end{quote}

Disorder-PDB has 233 targets and Disorder-NOX 178; both give 7. CAID2's
Disorder-PDB has 256 and also gives 7. **Collecting targets does not buy
resolution the labels lack.**


\subsection*{Proof outline}

The scores of a certified family must be pairwise separated by more than the
noise can move them, which is \code{c} in the units of the score. On a grid of
denominator \code{n}, values separated by more than \code{c} are separated by at least
\code{⌊c·n⌋ + 1} grid steps, and \code{[0, 1]} contains \code{n} steps. Counting the maximum
number of points one can place on \code{n} steps at that spacing gives the bound; the
attainment instance places them exactly at the spacing.


\subsection*{The converse, which is what makes this impossibility rather than low power}

\code{unresolvable\_pair}. A benchmark observes the \textbf{annotation}, not the truth.
Every labelling within the noise budget of the annotation is a consistent
truth. For two methods whose measured scores differ by less than the budget, the
proof \emph{constructs two such truths} --- one under which method A is strictly better
and one under which B is --- both compatible with every observation the benchmark
recorded.

The ordering is therefore \textbf{not a function of the data}. This is the reason no
amount of resampling helps: a bootstrap, a permutation test, leave-one-out and
Monte Carlo all estimate the sampling distribution of a statistic computed from
the same observations, and the ordering is not determined by those observations.
More resamples shrink the Monte Carlo error in locating a p-value and leave the
p-value itself alone. \textbf{Only new data or better labels move this bound.}

\code{over\_capacity\_has\_close\_pair} closes the loop: once the entry list exceeds the
capacity, a pair inside the budget necessarily exists. With 117 entrants and a
capacity of 7, CAID3 is over capacity by a factor of seventeen.


\subsection*{Scope, stated rather than assumed}

\begin{itemize}\itemsep1pt \parskip0pt
\item The noise model is a \textbf{worst-case budget}, not a probability model. That is
  precisely why the converse is a hard impossibility rather than a power
  statement; a stochastic model would give a different and weaker theorem.
\item "At capacity" is the maximum size of a \textbf{certifiable family}. It is never
  promoted to a claim that such a family has been certified.
\item \code{ε = 0.0801} is measured on 160 of 319 Disorder-PDB targets --- those with a
  usable MobiDB record --- pooled over 61,013 residues. The 159 remaining targets
  are excluded from the estimate, not assumed noise-free.
\item \code{δ = 0} throughout. Any non-zero effect-size threshold only lowers capacity.
\item The theorem bounds what a benchmark can certify about **error-count or
  averaged-score** rankings. It does not bound a \emph{paired} ranking statistic,
  which is the loophole Note S2 exploits deliberately and openly.
\end{itemize}


\vspace{0.5em}\hrule\vspace{0.8em}



\section*{Note S2 --- The second-order identity}


\subsection*{Statement}

Let \code{T ⊆ ι} be the truth on one protein (the disordered residues) and \code{L ⊆ ι}
the annotation. A pair \code{(p, q)} with \code{p ∈ T}, \code{q ∉ T} is \textbf{discordant} if the
annotation orders it the other way: \code{p ∉ L} and \code{q ∈ L}. Write \code{d = |T \ L|}
(down-flips: truly disordered, annotated ordered) and \code{u = |L \ T|} (up-flips).
Then, with no hypotheses at all:

\begin{quote}\ttfamily\footnotesize\begin{flushleft}
discordant(T,~L)~=~2~·~d~·~u~~~~~~~~`discordant\_eq\_flip\_product`\\
noise(T,~L)~~~~~~=~d~+~u~~~~~~~~~~~~`noise\_eq\_flip\_add`\\
discordant~~~~~~~≤~noise²~/~2~~~~~~~`discordant\_le\_noise\_sq`\\
\end{flushleft}\end{quote}

with equality in the third exactly when the flip classes balance
(\code{discordant\_eq\_of\_balanced}, \code{d = u}); \code{sharp\_instance} realises it on four
residues.

The pairs a pairwise protocol can score are the ones \textbf{both} labellings order,
and they are counted exactly too:

\begin{quote}\ttfamily\footnotesize\begin{flushleft}
|comparable(T,~L)|~=~2~·~(|T~∩~L|~·~|(T~∪~L)ᶜ|~+~d~·~u)~~~`card\_comparablePairs`\\
\end{flushleft}\end{quote}

so the pair noise rate is exactly

\begin{quote}\ttfamily\footnotesize\begin{flushleft}
ν\_pair~=~d·u~/~(|T~∩~L|~·~|(T~∪~L)ᶜ|~+~d·u).\\
\end{flushleft}\end{quote}

\code{noise\_orderKey} shows this quantity \textbf{is} the label noise of the answer key
the annotation induces on those pairs, which is what lets the residue-level
capacity machinery apply at the pair level verbatim: \code{pairwise\_capacity\_bound}
gives a certified family of at most \code{max 1 ⌈1/(2 ν\_pair)⌉} members, with no
hypotheses.


\subsection*{The closed form, and its hypotheses}

\begin{quote}\ttfamily\footnotesize\begin{flushleft}
nuPair\_le\_two\_eps\_sq~:~~~~~~~~ν\_pair~≤~2~ε²\\
pairwise\_capacity\_quadratic~:~ν\_pair~≤~2~ε²~~∧~~⌈1/(4~ε²)⌉~≤~pairwiseCapacity\\
\end{flushleft}\end{quote}

both require \code{BalancedClasses T L} --- among the residues the two labellings agree
on, as many are disordered as ordered, \code{|T ∩ L| = |(T ∪ L)ᶜ|} --- and \code{ε ≤ 1/4};
the capacity half also needs both flip classes nonempty. Balance is
load-bearing, not decoration: the denominator \code{|T ∩ L|·|(T ∪ L)ᶜ|} collapses
when the classes are lopsided and the gain collapses with it.
\code{capacity\_example} exhibits sixteen residues where every hypothesis holds and
the capacity is 25 against a residue capacity of 4, so the statement is not
vacuous.


\subsection*{Measurement}

MobiDB publishes the missing-residue call of each individual deposited structure
(\code{derived-missing\_residues-mobi-\{PDBID\}\_\{CHAIN\}}). For each protein we enumerate
all pairs of its structures and restrict to residues both cover: 159 CAID3
Disorder-PDB targets, 2,746 structure pairs.

\begin{center}\footnotesize
\begin{tabular}{lrr}
\toprule
\textbf{} & \textbf{pooled} & \textbf{median protein} \\
\midrule
\code{ε\_label} & 0.0651 & 0.0342 \\
\code{ε\_pair} & \textbf{0.0100} & 0.0006 \\
ratio & 6.54 & 57 \\
capacity \code{⌈1/2ε⌉} & \textbf{8 $\rightarrow$ 51} &  \\
\bottomrule
\end{tabular}
\end{center}

Per-protein values for all 159 targets are in \textbf{Table S7}.

\textbf{The hypotheses of the closed form were checked, not assumed, and they fail.}

\begin{center}\footnotesize
\begin{tabular}{lrr}
\toprule
\textbf{} & \textbf{structure pairs} & \textbf{} \\
\midrule
checked & 2,746 &  \\
balanced agreement classes within 10\% & \textbf{54} & 2.0\% \\
noise rate \code{ε ≤ 1/4} & 2,606 & 94.9\% \\
\textbf{both hypotheses} & \textbf{32} & \textbf{1.2\%} \\
bound \code{ν\_pair ≤ 2ε²} holds & 2,248 & 81.9\% \\
\textbf{\dots among the 32 satisfying both} & \textbf{32} & \textbf{100\%} \\
\bottomrule
\end{tabular}
\end{center}

Pooled, the agreement classes stand at 240,506 disordered against 555,402
ordered --- a ratio of 0.433 on a reference that is 31.6\% disordered --- and the
measured rate exceeds the closed form, 0.00996 against 0.00848. The bound holds
wherever the theorem says it holds and nowhere is it violated under its own
hypotheses. What CAID3 cannot use is the shortcut: **the gain has to be measured
on each benchmark rather than predicted from its label-noise rate.**


\subsection*{Two things this got wrong, and what caught each}

\textbf{The mechanism.} An earlier version attributed the gap to spatial correlation:
disorder flips in blocks, so only pairs crossing a block boundary reverse, and
the reduction should scale as \code{1/L}. The block statistics are real --- 7,100
disagreement blocks across 147 proteins, median length 3, mass-weighted mean
73.8, 67.4\% of context-dependent residues in blocks of ten or more --- but
\code{1/73.8 = 0.0136} against a measured ratio of \textbf{0.1075} over-predicts the
reduction eightfold. The identity above needs no correlation assumption.

\textbf{The number.} An earlier measurement put \code{ε\_pair} at 0.0070 and the capacity at
72, on a denominator that counted the pairs ordered by \emph{either} labelling with
the discordant ones counted twice,

\begin{quote}\ttfamily\footnotesize\begin{flushleft}
2ae~+~4du~+~(a+e)(d+u)~~~instead~of~~~2(ae~+~du),\\
\end{flushleft}\end{quote}

which is larger than the comparable set and mixes ordered with unordered counts.
Formalising \code{card\_comparablePairs} is what exposed it: the two expressions are
not equal. On the correct denominator the rate is \textbf{1.42$\times$ larger} and the
capacity is \textbf{51}. Both numbers are kept in \code{relative\_noise.json}
(\code{eps\_pairwise} and \code{eps\_pairwise\_superseded}) so the correction is auditable.


\section*{Note S3 --- The protocol, worked}


\subsection*{Specification}

As in the main text, steps 1--6. Two points deserve expansion.

\textbf{Why the unweighted mean (step 3).} \code{AUC\_within} as it appears in the
decomposition identity is the \textbf{pair-weighted} mean of per-target AUCs ---
weighted by each target's own positive$\times$negative pair count --- because the
identity does not close otherwise. Pair weighting lets a handful of long chains
carry the number. The protocol therefore scores the \textbf{unweighted} mean: one
protein, one vote. Both are calibration-invariant, since the invariance holds
target by target before any averaging; the relevant theorem is the per-target
form

\begin{quote}\ttfamily\footnotesize\begin{flushleft}
theorem~auc\_target\_strictMono\_invariant~(f~:~ℝ~→~ℝ)~(hf~:~StrictMono~f)~:\\
~~~~aucOn~t~(f~∘~s)~=~aucOn~t~s\\
\end{flushleft}\end{quote}

from which the weighted and unweighted means both follow. The published theorem
\code{auc\_within\_strictMono\_invariant} is stated for the pair-weighted form, so the
per-target statement is what the protocol should cite and it is pending
formalisation. This distinction is recorded rather than elided.

\textbf{Why coverage is a gate and not a covariate (step 1).} Declining a target
raises a within-protein score, because the targets methods decline are harder
(Table S9). A method that answers 87\% of Linker is not scored on Linker; it is
scored on an easier benchmark with the same name.


\subsection*{Worked example: what the protocol refuses to order}

Disorder-PDB, 120 entered, 60 eligible, 233 two-class targets, $\varepsilon$\_pair = 0.0100,
capacity $\lceil$1/(2$\cdot$0.0100)$\rceil$ = 51.

The leader, DisorderNet-pbias at 0.9512, is separated at the Holm-corrected
Wilcoxon margin from \textbf{55 of the 59} other eligible methods. The four it is not
separated from are reported as tied with it at rank 1 --- not as ranks 2, 3, 4 ---
and that is the whole of the protocol's honesty mechanism. The residue-level
protocol, on the same data, can certify \textbf{seven} orderings in total.


\subsection*{Rank changes, five references}

Full table in \textbf{Table S3}. The largest single movements:

\begin{center}\footnotesize
\begin{tabular}{llrrr}
\toprule
\textbf{method} & \textbf{benchmark} & \textbf{pooled} & \textbf{pairwise} & \textbf{move} \\
\midrule
LIPNet & Binding & 60 & 6 & \textbf{+54} \\
ProBiPred-nucleic & Binding & 62 & 17 & +45 \\
DisorderUnetLM & Binding-IDR & 67 & 23 & +44 \\
LINKER-Pred2 & Binding-IDR & 29 & 66 & \textbf{−37} \\
rawMSA & Binding & 4 & 40 & −36 \\
flDPnn2 & Binding-IDR & 64 & 29 & +35 \\
UdonPred-TriZOD & Linker & 44 & 11 & +33 \\
\textbf{AlphaFold-pLDDT} & \textbf{Disorder-NOX} & \textbf{39} & \textbf{6} & \textbf{+33} \\
\bottomrule
\end{tabular}
\end{center}


\vspace{0.5em}\hrule\vspace{0.8em}



\section*{Note S4 --- Pre-registered variants and their outcomes}

Every architectural idea in this project was registered before its run, with
floors, endpoints and a decision rule fixed in advance. \textbf{Five were rejected.}
They are reported here in full because the capacity result explains them, and
because a paper that reports only the variants that worked is the failure mode
this project was built to avoid.

\begin{center}\footnotesize
\begin{tabular}{rlllrl}
\toprule
\textbf{\#} & \textbf{variant} & \textbf{registered in} & \textbf{primary endpoint, as registered} & \textbf{measured} & \textbf{outcome} \\
\midrule
1 & \textbf{Private trunk} --- detach the binding tasks from the shared trunk & PREREG\_4 & Binding-IDR pooled $\geq$ 0.5900 & 0.5800 & \textbf{REJECTED}, plus three floors failed \\
2 & \textbf{Backbone handedness} --- two chirality channels from C$\alpha$ virtual torsions & PREREG\_5, \_6 & A/B against a regime-matched control on the fixed holdout & null on every task & \textbf{REJECTED} --- the gate opened 22$\times$ over threshold and the A/B saw nothing \\
3 & \textbf{Motif-scale read-out} --- a narrow private stack for the binding tasks & PREREG\_7 & beat \code{mt\_control} on pooled Binding-IDR & 0.5281 vs 0.5422 (−0.0141) & \textbf{REJECTED} \\
4 & \textbf{Distribution matching} --- a W₁ term on the protein-level score distribution & PREREG\_8 & beat \code{mt\_control} on \textbf{between-protein} Disorder-NOX & 0.8564 vs 0.8600 (−0.0036) & \textbf{REJECTED} \\
5 & \textbf{Within-protein ranking loss} --- optimise the calibration-invariant part directly & PREREG\_9 & beat \code{mt\_control} on \textbf{within-protein} Disorder-NOX & 0.8500 vs 0.8511 (−0.0011) & \textbf{REJECTED} \\
6 & \textbf{Ensemble} --- equal-weight rank fusion, membership chosen on the holdout & PREREG\_10 & beat PUNCH2 on pooled Disorder-PDB & +0.0062, p = 0.120 & \textbf{not significant} \\
\bottomrule
\end{tabular}
\end{center}

Endpoints 3--5 were computed for this paper (job 30098857,
\code{results/caid3/registered\_endpoints.py}), because each registration states its
endpoint on a statistic the CAID3 evaluator does not write and each had been
filed as rejected on the pooled score instead. **All three fail on the statistic
they were registered on, and not one of the fifteen paired per-target
comparisons against the control survives Holm** (smallest adjusted p = 0.345).

The two mechanism predictions are the sharper part. PREREG\_8 said the
distribution-matching term should move the \textbf{between}-protein component,
because the term is permutation-invariant within a protein and therefore cannot
touch the within-protein ordering. PREREG\_9 said the ranking surrogate should
move the \textbf{within}-protein component, tested directly rather than inferred from
a pooled number. Each registration named the quantity its own mechanism had to
move; neither moved it.

\textbf{The benchmark's verdict, checked against new data.} The capacity result says
CAID3 could not have separated these variants whatever they did, which invites
the objection that real improvements are being hidden. New data is not bounded by
CAID3's capacity, so the four scorable variants were scored on the 186 temporally
held-out chains (job 30102022), against the same regime-matched control:

\begin{center}\footnotesize
\begin{tabular}{lrrr}
\toprule
\textbf{variant vs \code{mt\_control}, 186 unseen chains} & \textbf{pooled} & \textbf{within (pair-wtd)} & \textbf{within (unwtd)} \\
\midrule
\code{mt\_motif} & −0.0033 & −0.0048 & −0.0075 \\
\code{mt\_rank} & −0.0035 & −0.0047 & −0.0137 \\
\code{mt\_wass} & −0.0046 & −0.0098 & −0.0084 \\
\code{mt\_chiral} & −0.0086 & −0.0301 & −0.0135 \\
\bottomrule
\end{tabular}
\end{center}

**Not one of them beats the control, on any of the three statistics, on data
none of them saw.** There was nothing for the benchmark to resolve. This is the
outcome the capacity theorem should be read as licensing: it says when to stop
looking, and when we checked with independent data, stopping was correct.

\textbf{An unregistered finding from the same computation.} On Binding-IDR,
restricting to the 42 targets that carry both classes among evaluated residues
reverses the sign of all three comparisons: \code{mt\_motif} −0.0141 $\rightarrow$ \textbf{+0.0480},
\code{mt\_wass} −0.0308 $\rightarrow$ \textbf{+0.0535}, \code{mt\_rank} −0.0050 $\rightarrow$ \textbf{+0.0117}. Ten of the 52
targets carry one class, and a single-class target contributes no within-protein
pair at all --- every pair it enters is between-protein. The whole of each
variant's apparent regression lives in targets that carry no residue-level
question. The registered endpoints still fail, because they were registered on
the pooled statistic and that is the statistic that decides them; but the reason
they fail is not that the variants are worse at the residue-level question. This
is the paper's thesis appearing where nobody put it.

The chirality result is the most instructive. The registered gate --- a
training-free handedness signal must beat an achiral control by $\geq$0.010
within-protein --- opened by a factor of twenty-two (+0.2256 on Disorder-PDB,
+0.2529 on Disorder-NOX), with the fitted direction the one polymer physics
predicts: disordered residues less right-handed (mean 0.074) than ordered
(0.211), and 61.3\% of torsions right-handed overall, as L-amino-acid proteins
should give. **And the A/B against a regime-matched control was null on every
task** (Disorder-PDB +0.0028, p = 0.087; Disorder-NOX +0.0019, p = 0.893;
Binding +0.0496, p = 0.200). The signal is real and the backbone already
carries it.

**Five architectures and an ensemble, each pre-registered, all failing to
separate from PUNCH2 on pooled AUC** (+0.0043, +0.0019, +0.0035, +0.0063,
+0.0062) is exactly what a benchmark over its capacity produces. We were trying
to resolve a pair the benchmark cannot resolve. On the calibration-invariant
axis the same comparison separates on all five references (main text, Fig. 4a).


\vspace{0.5em}\hrule\vspace{0.8em}



\section*{Note S5 --- Errors found and corrected}

Kept in full, because they are the reason to believe the rest, and because each
was caught by a measurement rather than by review.

\begin{center}\footnotesize
\begin{tabular}{lll}
\toprule
\textbf{what looked true} & \textbf{what was true} & \textbf{what caught it} \\
\midrule
Ten CAID3 numbers & the evaluator built the head with default dilations while the checkpoints were trained WIDE; dilation changes no weight shape, so \code{load\_state\_dict(strict=True)} accepted it silently and computed a different function at under a third of the trained context & re-deriving the receptive field from the checkpoint's own recorded flags \\
Homologues removed from the validation holdout & \code{homologous = set(hits)} was a set of \code{(query\_index, subject\_index, identity)} triples, so \code{id not in homologous} was always true; not one homologue left training while the log printed a plausible count & a log line reporting 1,157 held out against 2,538 reported homologues \\
"the head has no channel for protein-level information; the receptive field is 213 residues" & GroupNorm pools over the length axis, so the dependency span is \textbf{global}; perturbing residue 0 moves the logit at residue 400 & \code{lite\_head.measured\_dependency\_span}, written to test the claim \\
pairwise discordance is low because noise is block-correlated, \code{≈1/L} & the mechanism is second-order, \code{2·d·u}; \code{1/L} over-predicts the reduction eightfold & measuring both the block distribution and the ratio \\
\code{ε\_pair = 0.0070}, pairwise capacity 72 & the denominator counted the pairs ordered by \emph{either} labelling with the discordant ones counted twice. \code{card\_comparablePairs} counts it exactly; the rate is 1.42$\times$ larger and the capacity is \textbf{51} & formalising the denominator \\
\code{ε\_pair ≤ 2·ε\_label²}, "the bound holds and is tight to 20\%" & the bound needs the \emph{agreement} classes to balance, which holds on 32 of 2,746 structure pairs; pooled, the measured rate exceeds it. It holds on all 32 that qualify & formalising the hypotheses, then checking them \\
AlphaFold-rsa's disordered-residue coverage collapses to 0.527 and it fails the risk guarantee & artefact of scoring our models on 186 chains and the baselines on 61, with different splits; on identical chains rsa covers 0.994 and achieves the guarantee & matching the chain sets \\
every Linker method must flag 100\% of the protein & \code{α} below \code{1/(n+1)}: arithmetic about \emph{n}, not a fact about any method. Now an explicit refusal by name & the uniformity of the result \\
a dihedral implementation passes its mirror test & it was wrong twice --- a 180° offset, then a sign inversion --- and the mirror test passed both times & an independently written second implementation \\
a source-text test passes, so the code path works & \code{LEADERS[task]} is a 3-tuple unpacked as a pair; the test asserted on source text and the code raised at runtime & running it \\
p = 0.0 between two identical methods & both bootstrap tails must carry the ties & comparing a method with itself \\
a certified irreducible-error ranking & \code{crossed\_bound\_loose\_unbounded}: the greedy crossed-matching certificate is loose by an unbounded factor, because the obstruction is cyclic across three proteins and the certificate is pairwise & the Lean development, which proved the negative we did not ask for \\
"we have never beaten PUNCH2 significantly", stated four times & \code{mt\_pbias} beats PUNCH2 on pooled Disorder-PDB by +0.0083, 95\% CI [+0.0010, +0.0176], p = 0.0226. The checkpoint had been filed as rejected by its own pre-registration on three other benchmarks, so its Disorder-PDB comparison stopped being quoted at all & re-reading every checkpoint's result file \\
\bottomrule
\end{tabular}
\end{center}


\vspace{0.5em}\hrule\vspace{0.8em}



\section*{Note S6 --- Inventory of the formal development}

Lean 4, sorry-free, axioms \code{propext}, \code{Classical.choice}, \code{Quot.sound}
throughout.

\textbf{Metric decomposition and invariance}
\code{auc\_pooled\_decomp} $\cdot$ \code{auc\_within\_strictMono\_invariant} $\cdot$
\code{auc\_pooled\_shift\_diff} $\cdot$ \code{inversion\_requires\_between\_gap} $\cdot$
\code{auc\_shift\_le\_ceiling} $\cdot$ \code{ceiling\_gap\_of\_crossed\_matching} $\cdot$
\code{exists\_order\_ge} $\cdot$ \code{exists\_bias\_of\_order} $\cdot$ \code{exists\_gapMatrix} $\cdot$
\code{feasible\_iff\_gap} $\cdot$ \code{Example.inversion} (worked instance, reproduced exactly by
the analysis code)

\textbf{Label noise and certification}
\code{LabelNoise.ranking\_certified} $\cdot$ \code{RobustCertificate.certificate\_under\_mean\_error}

\textbf{Benchmark capacity}
\code{BenchmarkCapacity.card\_le\_benchCapacity} (grid / integer / real forms) $\cdot$
\code{benchCapacity\_attained} (three attainment instances) $\cdot$
\code{benchCapacity\_noise\_only} $\cdot$ \code{card\_le\_capacity\_labelNoise} $\cdot$
\code{unresolvable\_pair} $\cdot$ \code{over\_capacity\_has\_close\_pair}

\textbf{Pairwise scoring} (\code{DiscordantPairs.lean})
\code{discordant\_eq\_flip\_product} $\cdot$ \code{noise\_eq\_flip\_add} $\cdot$ \code{discordant\_le\_noise\_sq}
(and \code{\_real}) $\cdot$ \code{discordant\_eq\_of\_balanced} $\cdot$ \code{sharp\_instance} $\cdot$
\code{comparablePairs} / \code{orderKey} / \code{noise\_orderKey} $\cdot$ \code{card\_comparablePairs} $\cdot$
\code{card\_split} $\cdot$ \code{pairwise\_capacity\_bound} $\cdot$ \code{BalancedClasses} $\cdot$
\code{nuPair\_le\_two\_eps\_sq} $\cdot$ \code{pairwise\_capacity\_quadratic} $\cdot$
\code{capacity\_at\_eps\_0651} $\cdot$ \code{exampleHypotheses} / \code{exampleCounts} /
\code{capacity\_example}

\textbf{Multiple testing under dependence} (\code{DependentBH.lean},
\code{DependentBHSharp.lean}, \code{RegionScreen.lean})
\code{wgt\_decomp} $\cdot$ \code{fdp\_eq\_sum\_wgt} $\cdot$ \code{SelfConsistentP} / \code{bhList\_selfConsistent} $\cdot$
\code{mean\_wgt\_le} $\cdot$ \code{selfConsistent\_fdr\_le\_harmonic} $\cdot$ \code{benjamini\_yekutieli} /
\code{benjamini\_yekutieli\_level} $\cdot$ \code{bh\_uncorrected\_bound} $\cdot$
\code{by\_dominates\_bonferroni} $\cdot$ \code{by\_fdp\_tail} $\cdot$ \code{harm\_le\_one\_add\_log} /
\code{log\_le\_harm} / \code{harm\_thousand\_le\_eight} $\cdot$ \code{superuniform\_grid} /
\code{not\_superuniform\_of\_const} $\cdot$ \code{card\_windowsThrough} / \code{pval\_superuniform} /
\code{bhR\_pv} / \code{bhRej\_pv} / \code{bh\_fdr\_eq\_harmonic} $\cdot$ \code{harmonic\_factor\_sharp} $\cdot$
\code{uncorrected\_bh\_exceeds\_level} $\cdot$ \code{by\_level\_is\_attained} $\cdot$
\code{hypotheses\_satisfiable} $\cdot$ \code{Blocks} / \code{lift} / \code{card\_lift} / \code{lift\_inter} /
\code{stdBlocks} $\cdot$ \code{fdp\_lift} $\cdot$ \code{region\_screen\_fdr\_le} $\cdot$
\code{region\_level\_is\_less\_conservative} $\cdot$ \code{harmonic\_gain\_log}

\textbf{Calibration and prediction sets}
\code{Calibration.risk\_decomposition} $\cdot$ \code{risk\_eq\_resolution\_of\_calibrated} $\cdot$
\code{PredictionSets.validity\_is\_free}

\textbf{Complexity}
\code{Complexity/biasThreshold\_hard} --- NP-hardness of the recalibration optimum,
proved from a verifier definition of NP through circuit satisfiability, a
verified Tseitin transformation, independent set and weighted linear ordering,
each link a total reduction with a machine-checked correctness proof and a
machine-checked polynomial size bound. Lean has no cost model; polynomiality is
carried by those size bounds on structurally simple, explicitly given functions
rather than proved as running time, and this is documented in the development.
\code{AUCCrossedMatchingLooseness.crossed\_bound\_loose\_unbounded} --- the companion
negative.

\textbf{Distributional design} (used by PREREG\_8, and for the ensemble paper)
\code{DistributionVerdict.distributional\_design\_law} $\cdot$
\code{transportCost\_line\_eq\_cdfL1} $\cdot$ \code{DistanceRealizability} $\cdot$
\code{DistanceCutEnsembles}

\textbf{Pending, and used in the text ahead of formalisation}
\code{auc\_target\_strictMono\_invariant} --- the per-target form of the AUC invariance,
which is what the protocol's step 3 needs; the published theorem is stated for
the pair-weighted mean, and the invariance holds target by target before any
averaging. $\cdot$ \code{crc\_valid} --- conformal risk control, cited rather than
formalised.

The three items this list carried in an earlier draft ---
\code{discordant\_eq\_flip\_product}, \code{discordant\_le\_noise\_sq} and
\code{pairwise\_capacity\_quadratic} --- are now proved, and proving them changed two of
this paper's numbers (Note S2).


\vspace{0.5em}\hrule\vspace{0.8em}



\section*{Note S7 --- The screen: arbitrary dependence, and the unit of testing}


\subsection*{What was missing}

Two earlier parts of this development control the false discovery rate of a
proteome-scale screen for disordered regions, and each left the same hole. One
proves the Benjamini--Hochberg theorem in a \textbf{product} experiment --- one
independent coordinate per candidate. The other drops independence entirely but
changes the currency: each candidate must supply an \textbf{e-value}, not a p-value.
A real screen reports p-values, and its candidates share a calibration run, a
model and a training set, so the dependence is neither absent nor
sign-constrained. Both files recorded the gap in the same words: *under
arbitrary dependence the Benjamini--Hochberg procedure needs the harmonic
correction, which is not proved here.*


\subsection*{The theorem}

Proved in the same finite, measure-theory-free setting (\code{Law Ω}: an arbitrary
probability distribution on a finite outcome space), assuming only
\code{Superuniform P X} --- \code{P(X ≤ t) ≤ t} for every \code{t ≥ 0}, the validity of a **true
null's** p-value. Nothing is assumed about the non-nulls, the number of
candidates, or the joint law.

\begin{itemize}\itemsep1pt \parskip0pt
\item \code{selfConsistent\_fdr\_le\_harmonic} --- any self-consistent step-up rule at level
  \code{q}: \code{E[FDP] ≤ q · H\_m · |H₀| / m}.
\item \code{benjamini\_yekutieli}, \code{benjamini\_yekutieli\_level} --- BH at the deflated level
  \code{α / H\_m} controls the rate at \code{α}, whatever the dependence.
\item \code{bh\_uncorrected\_bound} --- uncorrected BH at level \code{q}: \code{E[FDP] ≤ q · H\_m}.
\item \code{by\_dominates\_bonferroni} --- the corrected list still contains the Bonferroni
  list at the same overall level, so the correction never loses a discovery to
  the family-wise procedure.
\item \code{harm\_le\_one\_add\_log}, \code{log\_le\_harm} --- the price is
  \code{log(m+1) ≤ H\_m ≤ 1 + log m}: logarithmic in the size of the screen, where
  Bonferroni's is linear. \code{harm\_thousand\_le\_eight}: a factor under 8 for a
  thousand candidates, against Bonferroni's 1000.
\item \code{superuniform\_grid} / \code{not\_superuniform\_of\_const} --- the assumption is
  satisfiable and has content.
\end{itemize}

\textbf{The mechanism}, isolated as \code{wgt\_decomp}, is why the joint law never enters.
Writing \code{R(ω)} for the reported list, the false discovery proportion is the
total weight \code{Σ\_\{i ∈ H₀\} 1\{i ∈ R\}/|R|} (\code{fdp\_eq\_sum\_wgt}), and the difficulty is
that \code{1/|R|} couples candidate \code{i} to every other. The identity

\begin{quote}\ttfamily\footnotesize\begin{flushleft}
1\{i~∈~R\}/|R|~=~Σ\_\{j=1\}\textasciicircum{}\{m\}~(1/j~−~1/(j+1))·1\{i~∈~R,~|R|~≤~j\}\\
~~~~~~~~~~~~~+~(1/(m+1))·1\{i~∈~R\}\\
\end{flushleft}\end{quote}

holds on every single outcome, and the events on the right are \textbf{nested}, not
disjoint slices \code{\{|R| = k\}}. Self-consistency --- every reported candidate has
\code{p\_i ≤ |R|·q/m}, which the BH list satisfies (\code{bhList\_selfConsistent}) --- turns
each of them into a statement about candidate \code{i} alone,

\begin{quote}\ttfamily\footnotesize\begin{flushleft}
\{i~∈~R,~|R|~≤~j\}~⊆~\{p\_i~≤~j·q/m\},\\
\end{flushleft}\end{quote}

whose probability validity alone bounds by \code{j·q/m}. Summing the telescoping
coefficients gives exactly \code{H\_m} (\code{mean\_wgt\_le}). No property of the joint law is
used anywhere, which is why the result holds under arbitrary dependence.


\subsection*{The factor is necessary, not conservative}

For every screen size \code{m} and every level \code{q}, \code{DependentBHSharp.lean}
constructs the joint law that attains it. Outcomes are \code{none} together with the
pairs \code{(j, s)}; in outcome \code{(j, s)} the candidates of the cyclic window of length
\code{j+1} starting at \code{s} have p-value \code{(j+1)q/m} and everyone else has \code{1}, and the
outcome carries probability \code{q/((j+1)·m)}.

\begin{itemize}\itemsep1pt \parskip0pt
\item The construction is balanced: each candidate lies in exactly \code{j+1} of the \code{m}
  windows of length \code{j+1} (\code{card\_windowsThrough}), so it sees the value
  \code{(j+1)q/m} with probability exactly \code{q/m}, for every \code{j}. Every one of the \code{m}
  p-values is therefore valid (\code{pval\_superuniform}) and all \code{m} hypotheses are
  true nulls.
\item In outcome \code{(j, s)} the observed p-values are \code{j+1} copies of \code{(j+1)q/m},
  which is exactly the step-up fixed point: BH stops at \code{j+1} and rejects
  precisely the window (\code{bhR\_pv}, \code{bhRej\_pv}). Every discovery is false.
\item The discovery outcomes of window length \code{j+1} carry total mass \code{q/(j+1)}, so
  \code{E[FDP] = Σ\_j q/j = q·H\_m} (\code{bh\_fdr\_eq\_harmonic}).
\end{itemize}

Consequences: the harmonic factor is attained (\code{harmonic\_factor\_sharp});
uncorrected BH at level \code{q} strictly exceeds its nominal level from \textbf{two}
candidates on (\code{uncorrected\_bh\_exceeds\_level}); and the corrected procedure sits
exactly at its own bound on the same instance (\code{by\_level\_is\_attained}), so
\code{benjamini\_yekutieli\_level} cannot be improved either.
\code{hypotheses\_satisfiable} supplies an admissible level for every \code{m ≥ 2}.


\subsection*{Where the price should be paid}

The factor depends on one thing only --- how many hypotheses the screen states ---
so the unit of testing is a design parameter, and disorder is a property of a
region.

\begin{itemize}\itemsep1pt \parskip0pt
\item \code{Blocks n M b} partitions the \code{n} residues into \code{M} candidate regions of \code{b}
  residues; disjointness and common size are all that is used. \code{stdBlocks}
  supplies the concrete partition, so nothing is vacuous.
\item \code{fdp\_lift} --- \textbf{the invariance}: reading a region-level report at residue
  level multiplies both the numerator and the denominator of the false
  discovery proportion by the region length, and it cancels. The two screens
  report the \emph{same} FDP. This is exactly where region-level annotation matters:
  it is true because the truth and the report are both unions of whole regions.
\item \code{region\_screen\_fdr\_le} --- so the region-level Benjamini--Yekutieli procedure at
  level \code{α/H\_M} controls the \textbf{residue-level} false discovery rate at \code{α} under
  an arbitrary joint law.
\item \code{region\_level\_is\_less\_conservative}, \code{harmonic\_gain\_log} --- and it tests at a
  strictly more permissive threshold, with the saving bounded below by
  \code{H\_n − H\_M ≥ log b − 1}: one nat per factor in the region length.
\end{itemize}


\subsection*{Instantiated on this paper's references}

\code{results/caid3/region\_screen.py}, job 30101613. Candidate regions are the
maximal same-label runs of each reference --- the unit the annotation is actually
constant on --- with a run broken by any unevaluated residue, which states more
hypotheses rather than fewer. Full table in \textbf{Table S15}.

\begin{center}\footnotesize
\begin{tabular}{lrrrrrrrr}
\toprule
\textbf{reference} & \textbf{n} & \textbf{M} & \textbf{mean region} & \textbf{H\_n} & \textbf{H\_M} & \textbf{saving} & \textbf{$\geq$ log b − 1} & \textbf{threshold} \\
\midrule
CAID3 Disorder-PDB & 99,239 & 1,015 & 97.8 & 12.083 & 7.500 & 4.582 & 3.583 & $\times$1.61 \\
CAID3 Disorder-NOX & 99,977 & 576 & 173.6 & 12.090 & 6.934 & 5.156 & 4.157 & $\times$1.74 \\
CAID3 Binding & 28,263 & 165 & 171.3 & 10.827 & 5.686 & 5.140 & 4.143 & $\times$1.90 \\
CAID3 Binding-IDR & 7,741 & 157 & 49.3 & 9.532 & 5.637 & 3.895 & 2.898 & $\times$1.69 \\
CAID3 Linker & 20,498 & 105 & 195.2 & 10.505 & 5.236 & 5.269 & 4.274 & $\times$2.01 \\
CAID2 Disorder-PDB & 130,877 & 1,162 & 112.6 & 12.359 & 7.636 & 4.724 & 3.724 & $\times$1.62 \\
CAID2 Disorder-NOX & 160,802 & 623 & 258.1 & 12.565 & 7.013 & 5.553 & 4.553 & $\times$1.79 \\
CAID2 Binding & 67,169 & 219 & 306.7 & 11.692 & 5.969 & 5.724 & 4.726 & $\times$1.96 \\
CAID2 Linker & 37,150 & 124 & 299.6 & 11.100 & 5.402 & 5.698 & 4.702 & $\times$2.05 \\
\bottomrule
\end{tabular}
\end{center}

The proved floor \code{log b − 1} is met with room on every reference. A region-level
screen may test at a threshold 1.61$\times$ to 2.05$\times$ larger for the same residue-level
false discovery rate, and by \code{fdp\_lift} it gives up nothing to do so.


\subsection*{The screen, run}

\code{results/caid3/conformal\_screen.py} (job 30102082) runs the procedure on CAID3
Disorder-PDB: 304 usable chains, 98,938 evaluated residues, candidates are
consecutive blocks of \code{b} residues, p-values are split conformal against the
known-ordered blocks of 80\% of chains, 50 random splits, medians reported.

\begin{center}\footnotesize
\begin{tabular}{rrlrrr}
\toprule
\textbf{b} & \textbf{$\alpha$} & \textbf{procedure} & \textbf{residues reported} & \textbf{realised residue FDP} & \textbf{power} \\
\midrule
100 & 0.10 & \textbf{region BY} ($\alpha$/H\_M) & \textbf{2,942} & 0.024 & \textbf{0.568} \\
100 & 0.10 & residue BY ($\alpha$/H\_n) & 203 & 0.000 & 0.036 \\
25 & 0.10 & \textbf{region BY} & 604 & 0.000 & 0.122 \\
25 & 0.10 & residue BY & 201 & 0.000 & 0.027 \\
10 & 0.10 & \textbf{region BY} & 562 & 0.000 & 0.098 \\
10 & 0.10 & residue BY & 204 & 0.000 & 0.034 \\
\bottomrule
\end{tabular}
\end{center}

At 100-residue blocks the region-level screen reports \textbf{14.5$\times$ the residues} of
the residue-level screen at the same guaranteed rate, and 2.75--3.0$\times$ at 10 and 25.
Uncorrected region-level BH reports far more (4,225 at b=100) but realises a
false discovery proportion of 0.099 against a nominal 0.10 with no guarantee
under dependence --- which is what the correction is being paid for.
Region-level Bonferroni reports nothing at any setting.

\code{fdp\_lift} is checked rather than assumed: the region-level and residue-level
false discovery proportions of the same report agree to within 0.006 at b=10 and
0.077 at b=100. The gap is not zero because the truth is not exactly a union of
whole blocks in real data --- a block is called disordered by majority --- and the
size of the gap measures how far the theorem's hypothesis is from holding.

\textbf{A third constraint dominates both, and it is not in the theory.} A conformal
p-value cannot fall below \code{1/(n\_cal + 1)}, and a step-up rule at level \code{α/H\_m}
over \code{m} hypotheses applies thresholds as small as \code{α/(m·H\_m)}. A screen is
therefore \emph{arithmetically} incapable of a discovery unless

\begin{quote}\ttfamily\footnotesize\begin{flushleft}
n\_cal~+~1~~≥~~m~·~H\_m~/~α~,\\
\end{flushleft}\end{quote}

whatever the predictor does. The requirement is \textbf{linear} in the number of
hypotheses, where the harmonic factor is logarithmic, so the unit of testing
sets the calibration budget far more strongly than it sets the correction.
Measured here, the region-level screen needs 13$\times$ fewer calibration units at
b=10, 36$\times$ at b=25, 80$\times$ at b=50 and 178$\times$ at b=100 than the residue-level screen.
Every configuration is still short of the requirement --- by 23--60$\times$ for regions
and 38--77$\times$ for residues --- and discoveries occur anyway only because many blocks
tie at the p-value floor. **This is the constraint that binds first in practice,
and a screen should be designed against it before the harmonic factor.**


\subsection*{Scope}

This is multiplicity control and it is orthogonal to calibration.
\code{Superuniform} is a statement about the \emph{null law of the read-out}; if the
instrument's baseline rate is wrong, the p-values are not valid and every bound
above is void. The screen reported above is a demonstration on a public benchmark whose labels
are already known; \textbf{no biological discovery is claimed}, and its realised false
discovery proportion sits far below nominal, as arbitrary-dependence control
should.


\vspace{0.5em}\hrule\vspace{0.8em}



\section*{Table S9 --- Coverage as a confound}

What declining targets is worth, measured against a fully-covering yardstick.

\begin{center}\footnotesize
\begin{tabular}{llrrr}
\toprule
\textbf{benchmark} & \textbf{method} & \textbf{declined} & \textbf{median length of declined} & \textbf{difficulty gap (AUC)} \\
\midrule
Disorder-NOX & ESMDisPred-2PDB & 23 / 204 & 1,292 & \textbf{+0.118} \\
Linker & ESMDisPred-2PDB & 6 / 31 & 1,214 & \textbf{+0.125} \\
Disorder-PDB & SPOT-Disorder2 & 21 / 319 & 1,684 & +0.035 \\
\bottomrule
\end{tabular}
\end{center}

Median length across all CAID3 targets is 344. Long proteins are harder and they
are what gets declined. This is why the protocol makes coverage a gate (step 1)
and why every rank in this paper is reported twice --- among all scored entrants
and among full-coverage entrants.


\section*{Table S10 --- Thresholded error counts}

Each method's own binary call (column four of its \code{.caid} file), against the
reference, over every evaluated residue. 46 methods submitted a complete binary
column.

\textbf{Disorder-PDB} (99,239 evaluated residues)

\begin{center}\footnotesize
\begin{tabular}{rlrrr}
\toprule
\textbf{\#} & \textbf{method} & \textbf{errors} & \textbf{rate} & \textbf{vs \#1} \\
\midrule
1 & \textbf{DisorderNet-pbias} & \textbf{8,000} & 0.081 & --- \\
2 & PUNCH2 & 8,460 & 0.085 & +460 \\
3 & DisorderNet-windowed & 8,481 & 0.085 & +481 \\
4 & PUNCH2-Light & 8,716 & 0.088 & +716 \\
5 & SETH-0 & 9,977 & 0.101 & +1,977 \\
\bottomrule
\end{tabular}
\end{center}

\textbf{Disorder-NOX} (99,977 evaluated residues)

\begin{center}\footnotesize
\begin{tabular}{rlrrr}
\toprule
\textbf{\#} & \textbf{method} & \textbf{errors} & \textbf{rate} & \textbf{vs \#1} \\
\midrule
1 & \textbf{DisorderNet-windowed} & \textbf{18,084} & 0.181 & --- \\
2 & DisorderNet-pbias & 18,983 & 0.190 & +899 \\
3 & flDPnn2 & 20,803 & 0.208 & +2,719 \\
4 & flDPlr & 22,382 & 0.224 & +4,298 \\
5 & ESpritz-D & 22,426 & 0.224 & +4,342 \\
\bottomrule
\end{tabular}
\end{center}

\textbf{Thresholds are not comparable across methods.} Each entrant chose its own,
and a method optimising AUC rather than accuracy is penalised here through no
fault of its ranking. This is reported beside the AUC tables rather than instead
of them.


\section*{Table S11 --- Certification frontier}

\code{LabelNoise.ranking\_certified} turns a margin into a statement about which
method is better \emph{in truth}: a margin above \code{2·ε·n} certifies the ranking at
annotation error rate \code{ε}. Because the criterion is monotone in \code{ε}, the whole
answer is a frontier and no point estimate of \code{ε} is needed.

\begin{center}\footnotesize
\begin{tabular}{lrr}
\toprule
\textbf{assumed annotation error} & \textbf{certified vs 45 others, Disorder-PDB} & \textbf{Disorder-NOX} \\
\midrule
0.25\% & 43 & 45 \\
1\% & 41 & 44 \\
\textbf{2\%} & \textbf{39} & \textbf{43} \\
3\% & 34 & 34 \\
5\% & 21 & 9 \\
\textbf{8.01\% (measured)} & \textbf{15} & \textbf{4} \\
\bottomrule
\end{tabular}
\end{center}

At the measured rate, \textbf{30 of 45} comparisons on Disorder-PDB and \textbf{41 of 45}
on Disorder-NOX are inside the annotation noise --- not unresolved by this
analysis, but unresolvable.

The \textbf{breakdown rate} of a comparison, \code{margin/(2n)}, is the annotation error
rate that would have to be exceeded before that ranking could be an artefact.
For the closest comparison it is 0.232\% on Disorder-PDB (PUNCH2, 460 residues)
and 0.450\% on Disorder-NOX; median across the field, 4.78\% and 3.71\%.


\section*{Table S12 --- Data provenance}

Reference FASTAs, from
\code{https://caid.idpcentral.org/assets/sections/challenge/static/references/3/}:

\begin{center}\footnotesize
\begin{tabular}{llrrrr}
\toprule
\textbf{file} & \textbf{md5} & \textbf{targets} & \textbf{positives} & \textbf{negatives} & \textbf{unevaluated} \\
\midrule
\code{disorder\_pdb.fasta} & \code{6feaff35263e7fd4a3f03640c23786fe} & 319 & 31,401 & 67,838 & 61,183 \\
\code{disorder\_nox.fasta} & \code{e2056cccfe186ab9da8d1e15f1508bef} & 204 & 26,367 & 73,610 & 541 \\
\code{binding.fasta} & \code{3094881568be4105b0e9afc576a79558} & 52 & 2,991 & 25,272 & 0 \\
\code{binding\_idr.fasta} & \code{4331f14502505e5dea58fcef0a29b2b7} & 52 & 2,991 & 4,750 & 20,522 \\
\code{linker.fasta} & \code{603e3055d5bed0d5af421a81123b252f} & 31 & 1,379 & 19,119 & 0 \\
\bottomrule
\end{tabular}
\end{center}

Prediction archive: \code{predictions.zip}, 76,480,806 bytes, 117 \code{.caid} files,
aggregate md5 of \code{md5sum *.caid} = \code{cdef7ae6dc3dd6ec8055aacdb14a54a7}.

\textbf{Published leaders reproduced through our pipeline} --- the check that licenses
every other number:

\begin{center}\footnotesize
\begin{tabular}{llrr}
\toprule
\textbf{challenge} & \textbf{leader} & \textbf{computed} & \textbf{published} \\
\midrule
Disorder-PDB & PUNCH2 & 0.9552 & 0.955 \\
Disorder-NOX & ESMDisPred-2PDB & 0.8855 & 0.885 \\
Binding & DisoFLAG-PB & 0.7760 & 0.776 \\
Binding-IDR & bindEmbed21IDR-rawGeneral & 0.6407 & 0.641 \\
Linker & IPA-AF2-Linker & 0.8985 & 0.897 \\
\bottomrule
\end{tabular}
\end{center}

Use the dataset named \textbf{"CAID3 v3"}. The API also serves an entry named plainly
"CAID3" (185 proteins), which is a different, smaller set and does not reproduce
the published leaderboard.


\section*{Table S13 --- Job registry}

\begin{center}\footnotesize
\begin{tabular}{llrl}
\toprule
\textbf{analysis} & \textbf{script} & \textbf{job} & \textbf{wall} \\
\midrule
temporal holdout, build & \code{rockfish/build\_temporal\_holdout.py} & 30034161 & 1 min \\
temporal holdout, AlphaFold prefetch & --- & 30034251 & 15 s \\
temporal holdout, evaluation & \code{rockfish/slurm/eval\_temporal.sbatch} & 30034521 & 2 m 47 s \\
conformal guarantee & \code{rockfish/slurm/eval\_temporal.sbatch} & 30035126 & 2 m 17 s \\
operating cost, whole field & \code{results/caid3/operating\_cost.py} & 30037497 & --- \\
registered endpoints, PREREG\_7/\_8/\_9 & \code{results/caid3/registered\_endpoints.py} & 30098857 & < 5 min \\
pairwise noise, corrected denominator & \code{results/caid3/relative\_noise.py} & 30101641 & < 5 min \\
region screen instantiation & \code{results/caid3/region\_screen.py} & 30101613 & < 5 min \\
pairwise leaderboard, re-run at $\varepsilon$\_pair = 0.0100 & \code{results/caid3/pairwise\_leaderboard.py} & 30101721 & < 5 min \\
registered variants on the temporal holdout & \code{rockfish/eval\_temporal.py} & 30102022 & 14 min \\
the screen, run end to end & \code{results/caid3/conformal\_screen.py} & 30102082 & < 5 min \\
chirality A/B, treatment & \code{rockfish/train\_multitask.py --chiral} & 30022469 & 11 h 08 m \\
chirality A/B, control & \code{rockfish/train\_multitask.py} & 30022470 & 11 h 20 m \\
private trunk & \code{rockfish/train\_multitask.py --private-attached} & 29922832 & 18 h 24 m \\
private trunk, evaluation & \code{rockfish/eval\_caid3\_official.py} & 30020129 & --- \\
windowed training & \code{rockfish/train\_multitask.py} & 29833016 & --- \\
windowed evaluation & \code{rockfish/eval\_caid3\_official.py} & 29833018 & --- \\
\bottomrule
\end{tabular}
\end{center}

Raw outputs are the JSON files under \code{paper/figures/}, each written by the job
named here and read directly by \code{paper/make\_figures.py} and
\code{paper/make\_supplementary.py}.
